% GENERATED FILE. Edit canonical lessons, then run npm run generate:books.
% canonical-source-hash: fnv1a64:416f1488c020c550
% canonical-lessons: GU-C164-dabbo, GU-C164-khad, GU-C164-mari, GU-C164-rai, GU-C164-elci, GU-R164-review-159-164

\chapter{The spice tin}
\label{ch:gu-the-spice-tin}
\hlchaptermodality{164}

\begin{chapteropening}
\textbf{By the end of this chapter:} \emph{I can name the spices in the kitchen tin.}

Complete the last lesson of chapter 164: I can name the spices in the kitchen tin.
\end{chapteropening}

\section[\texorpdfstring{\gu{ડબ્બો} (ḍabbo)}{ડબ્બો (ḍabbo)}]{\gu{ડબ્બો} (ḍabbo) — tin, box}

\label{lesson:GU-C164-dabbo}



\begin{quote}
Before the new one: say the Gujarati for to serve, then the Gujarati for to share out.
\end{quote}

\subsection*{You'll want to know: \gu{ડબ્બો}}
\textbf{\gu{ડબ્બો}} — \emph{ḍabbo} — \textquotedblleft{}tin, box\textquotedblright{}.

\begin{quote}
\textbf{\gu{મસાલાનો} \gu{ડબ્બો} \gu{રસોડામાં} \gu{છે}.} — \emph{masālāno ḍabbo rasoḍāmā̃ che} — \textquotedblleft{}The spice tin is in the kitchen.\textquotedblright{}
\end{quote}

\begin{cousinweb}
The origin of \textbf{\gu{ડબ્બો}} (\emph{ḍabbo}) is not settled. CDIAL lists no Gujarati form of it, and EtymDB has no entry for it.
\end{cousinweb}

\begin{grammarlens}[title={What you've built}]
One more word for everyday life at A2.
\end{grammarlens}

\subsection*{Your turn}
\begin{itemize}
\raggedright
  \item \emph{Say it:} \emph{ḍabbo}
  \item \emph{Say it:} \emph{ḍabbo}, once more
  \item \emph{Say it:} \emph{masālāno ḍabbo rasoḍāmā̃ che} — The spice tin is in the kitchen.
  \item \emph{Recall:} say the Gujarati for to serve, then the Gujarati for to share out, then say \emph{ḍabbo} again
\end{itemize}

\begin{tcolorbox}[breakable,colback=teal!4,colframe=teal!35!black,title={Before you move on}]
What does \gu{ડબ્બો} mean? (\textquotedblleft{}Tin, box\textquotedblright{}.) Say it once more.
\end{tcolorbox}
\section[\texorpdfstring{\gu{ખાંડ} (khā̃ḍ)}{ખાંડ (khā̃ḍ)}]{\gu{ખાંડ} (khā̃ḍ) — sugar}

\label{lesson:GU-C164-khad}



\begin{quote}
Before the new one: say the Gujarati for to share out, then the Gujarati for tin.
\end{quote}

\subsection*{You'll want to know: \gu{ખાંડ}}
\textbf{\gu{ખાંડ}} — \emph{khā̃ḍ} — \textquotedblleft{}sugar\textquotedblright{}.

\begin{quote}
\textbf{\gu{ડબ્બામાં} \gu{ખાંડ} \gu{નથી}.} — \emph{ḍabbāmā̃ khā̃ḍ nathī} — \textquotedblleft{}There is no sugar in the tin.\textquotedblright{}
\end{quote}

A feminine noun.

\begin{cousinweb}
\textbf{\gu{ખાંડ}} (\emph{khā̃ḍ}) is a tadbhava. It comes from Sanskrit \emph{khaṇḍa}, \emph{khaṇḍu}, \textquotedblleft{}sugar; partially dried molasses\textquotedblright{}, and CDIAL glosses the Gujarati word \textquotedblleft{}coarse sugar\textquotedblright{} (CDIAL 3800).
\end{cousinweb}

\begin{grammarlens}[title={What you've built}]
One more word for everyday life at A2.
\end{grammarlens}

\subsection*{Your turn}
\begin{itemize}
\raggedright
  \item \emph{Say it:} \emph{khā̃ḍ}
  \item \emph{Say it:} \emph{khā̃ḍ}, once more
  \item \emph{Say it:} \emph{ḍabbāmā̃ khā̃ḍ nathī} — There is no sugar in the tin.
  \item \emph{Recall:} say the Gujarati for to share out, then the Gujarati for tin, then say \emph{khā̃ḍ} again
\end{itemize}

\begin{tcolorbox}[breakable,colback=teal!4,colframe=teal!35!black,title={Before you move on}]
What does \gu{ખાંડ} mean? (\textquotedblleft{}Sugar\textquotedblright{}.) Say it once more.
\end{tcolorbox}
\section[\texorpdfstring{\gu{મરી} (marī)}{મરી (marī)}]{\gu{મરી} (marī) — black pepper}

\label{lesson:GU-C164-mari}



\begin{quote}
Before the new one: say the Gujarati for tin, then the Gujarati for sugar.
\end{quote}

\subsection*{You'll want to know: \gu{મરી}}
\textbf{\gu{મરી}} — \emph{marī} — \textquotedblleft{}black pepper\textquotedblright{}.

\begin{quote}
\textbf{\gu{મરી} \gu{તીખી} \gu{છે}.} — \emph{marī tīkhī che} — \textquotedblleft{}Black pepper is hot.\textquotedblright{}
\end{quote}

\begin{cousinweb}
\textbf{\gu{મરી}} (\emph{marī}) is a tadbhava. It comes from Sanskrit \emph{marica}, \textquotedblleft{}peppercorn\textquotedblright{}, and CDIAL glosses the Gujarati word \textquotedblleft{}black pepper\textquotedblright{} (CDIAL 9875).
\end{cousinweb}

\begin{grammarlens}[title={What you've built}]
One more word for everyday life at A2.
\end{grammarlens}

\subsection*{Your turn}
\begin{itemize}
\raggedright
  \item \emph{Say it:} \emph{marī}
  \item \emph{Say it:} \emph{marī}, once more
  \item \emph{Say it:} \emph{marī tīkhī che} — Black pepper is hot.
  \item \emph{Recall:} say the Gujarati for tin, then the Gujarati for sugar, then say \emph{marī} again
\end{itemize}

\begin{tcolorbox}[breakable,colback=teal!4,colframe=teal!35!black,title={Before you move on}]
What does \gu{મરી} mean? (\textquotedblleft{}Black pepper\textquotedblright{}.) Say it once more.
\end{tcolorbox}
\section[\texorpdfstring{\gu{રાઈ} (rāī)}{રાઈ (rāī)}]{\gu{રાઈ} (rāī) — mustard seed}

\label{lesson:GU-C164-rai}



\begin{quote}
Before the new one: say the Gujarati for sugar, then the Gujarati for black pepper.
\end{quote}

\subsection*{You'll want to know: \gu{રાઈ}}
\textbf{\gu{રાઈ}} — \emph{rāī} — \textquotedblleft{}mustard seed\textquotedblright{}.

\begin{quote}
\textbf{\gu{રાઈ} \gu{નાની} \gu{અને} \gu{કાળી} \gu{છે}.} — \emph{rāī nānī ane kāḷī che} — \textquotedblleft{}Mustard seed is small and black.\textquotedblright{}
\end{quote}

\begin{cousinweb}
\textbf{\gu{રાઈ}} (\emph{rāī}) is a tadbhava, a word inherited from Sanskrit and reshaped by sound change. It comes from \emph{rājikā}, \textquotedblleft{}mustard\textquotedblright{}, and CDIAL glosses the Gujarati word \textquotedblleft{}mustard\textquotedblright{} (CDIAL 10688).
\end{cousinweb}

\begin{grammarlens}[title={What you've built}]
One more word for everyday life at A2.
\end{grammarlens}

\subsection*{Your turn}
\begin{itemize}
\raggedright
  \item \emph{Say it:} \emph{rāī}
  \item \emph{Say it:} \emph{rāī}, once more
  \item \emph{Say it:} \emph{rāī nānī ane kāḷī che} — Mustard seed is small and black.
  \item \emph{Recall:} say the Gujarati for sugar, then the Gujarati for black pepper, then say \emph{rāī} again
\end{itemize}

\begin{tcolorbox}[breakable,colback=teal!4,colframe=teal!35!black,title={Before you move on}]
What does \gu{રાઈ} mean? (\textquotedblleft{}Mustard seed\textquotedblright{}.) Say it once more.
\end{tcolorbox}
\section[\texorpdfstring{\gu{એલચી} (elcī)}{એલચી (elcī)}]{\gu{એલચી} (elcī) — cardamom}

\label{lesson:GU-C164-elci}



\begin{quote}
Before the new one: say the Gujarati for black pepper, then the Gujarati for mustard seed.
\end{quote}

\subsection*{You'll want to know: \gu{એલચી}}
\textbf{\gu{એલચી}} — \emph{elcī} — \textquotedblleft{}cardamom\textquotedblright{}.

\begin{quote}
\textbf{\gu{એલચી} \gu{મોંઘી} \gu{છે}.} — \emph{elcī moṅghī che} — \textquotedblleft{}Cardamom is expensive.\textquotedblright{}
\end{quote}

\begin{cousinweb}
The origin of \textbf{\gu{એલચી}} (\emph{elcī}) is not settled. CDIAL lists Sanskrit \emph{ēlā}, \textquotedblleft{}cardamom\textquotedblright{}, which it takes from Dravidian, but no Gujarati form under it (CDIAL 2522).
\end{cousinweb}

\begin{grammarlens}[title={What you've built}]
One more word for everyday life at A2.
\end{grammarlens}

\subsection*{Your turn}
\begin{itemize}
\raggedright
  \item \emph{Say it:} \emph{elcī}
  \item \emph{Say it:} \emph{elcī}, once more
  \item \emph{Say it:} \emph{elcī moṅghī che} — Cardamom is expensive.
  \item \emph{Recall:} say the Gujarati for black pepper, then the Gujarati for mustard seed, then say \emph{elcī} again
\end{itemize}

\begin{tcolorbox}[breakable,colback=teal!4,colframe=teal!35!black,title={Before you move on}]
What does \gu{એલચી} mean? (\textquotedblleft{}Cardamom\textquotedblright{}.) Say it once more.
\end{tcolorbox}
\section[(dialogue)]{Words from chapters 159–164 — a review}

\label{lesson:GU-R164-review-159-164}



\begin{quote}
Say the Gujarati for to scour and to dry out.
\end{quote}

\subsection*{Your turn}
Say these aloud:

\begin{itemize}
\raggedright
  \item to scour, to dry out, water tap, to spill, floor
  \item bedding, to spread out, swing, electric light, to put out
  \item to move, to shake out, to cover, to break, to gather
  \item housework, bell, to tidy up, clay pot, corner
  \item breakfast, meal, to have a meal, to serve, to share out
  \item tin, sugar, black pepper, mustard seed, cardamom
\end{itemize}

\begin{tcolorbox}[breakable,colback=teal!4,colframe=teal!35!black,title={Before you move on}]
To scour? (\textbf{\gu{માંજવું}}.) To dry out? (\textbf{\gu{સૂકવવું}}.) Water tap? (\textbf{\gu{નળ}}.) To spill? (\textbf{\gu{ઢોળવું}}.) Floor? (\textbf{\gu{ભોંય}}.) Bedding? (\textbf{\gu{પથારી}}.) To spread out? (\textbf{\gu{પાથરવું}}.) Swing? (\textbf{\gu{હિંચકો}}.) Electric light? (\textbf{\gu{બત્તી}}.) To put out? (\textbf{\gu{હોલવવું}}.) To move? (\textbf{\gu{ખસેડવું}}.) To shake out? (\textbf{\gu{ખંખેરવું}}.) To cover? (\textbf{\gu{ઢાંકવું}}.) To break? (\textbf{\gu{તૂટવું}}.) To gather? (\textbf{\gu{ભેગું} \gu{કરવું}}.) Housework? (\textbf{\gu{ઘરકામ}}.) Bell? (\textbf{\gu{ઘંટડી}}.) To tidy up? (\textbf{\gu{સમેટવું}}.) Clay pot? (\textbf{\gu{માટલું}}.) Corner? (\textbf{\gu{ખૂણો}}.) Breakfast? (\textbf{\gu{નાસ્તો}}.) Meal? (\textbf{\gu{ભોજન}}.) To have a meal? (\textbf{\gu{જમવું}}.) To serve? (\textbf{\gu{પીરસવું}}.) To share out? (\textbf{\gu{વહેંચવું}}.) Tin? (\textbf{\gu{ડબ્બો}}.) Sugar? (\textbf{\gu{ખાંડ}}.) Black pepper? (\textbf{\gu{મરી}}.) Mustard seed? (\textbf{\gu{રાઈ}}.) Cardamom? (\textbf{\gu{એલચી}}.)
\end{tcolorbox}
